%----------------------------------------------------------------------
\section{Operation and exceptional conditions}
%----------------------------------------------------------------------
\textbf{(i) Missing observations.}
\begin{itemize}[leftmargin=*,itemsep=1pt]
\item Every value carries a mask bit and an age (Table~\ref{tab:data}). Missing levels are filled
with the last valid value and missing rain with zero, so the network can see that the value is
stale and is not misled by the fill.
\item Training randomly drops single readings, whole hours and entire stations (10--20\,\% of
samples), so the model learns to rely on neighbours, rainfall and the latent node when a sensor is
silent. The gate in block (3) learns when the station's own history cannot be trusted.
\item The forecast anchor $y_i(\tau^{\ast}_i)$ is the last valid reading. If a station has been
silent for more than one hour, its forecast is still issued from upstream information but is
flagged \emph{degraded}, and its quantile interval is wider (learned from the masked training
samples).
\item Missing targets are excluded from the loss ($M_{b,i,h}$) and never imputed for training.
\end{itemize}

\textbf{(ii) Delayed observations.}
\begin{itemize}[leftmargin=*,itemsep=1pt]
\item Inputs are always assembled \emph{as of} $\tau$ (A2): a reading that has not arrived is
treated as missing, with its age, and the system does not wait for it.
\item When the late reading arrives, the next issuance ten minutes later uses it automatically,
because the history window is rebuilt every cycle. Published forecasts are not rewritten. Each is
archived with a snapshot of the data available when it was made, so later evaluation reflects what
the system actually knew.
\item Training replays the empirical delay distribution, so the model is trained with the same
information it will have in operation. Training on the complete, cleaned archive would overstate
accuracy.
\end{itemize}

\textbf{(iii) Announced future releases.}
\begin{itemize}[leftmargin=*,itemsep=1pt]
\item The latest schedule announced at or before $\tau$ is placed on the future grid
$\tau+\Delta..\tau+6\Delta$ by \emph{effective} time (as-of join, Section~3) and enters through the
reservoir's future extension. It is the only input from after $\tau$.
\item A ``time since announcement'' feature lets the model learn that recent announcements are
more reliable than old ones.
\item If no schedule covers part of the horizon, the last actual release is held (A5) and a mask
bit marks the gap.
\item A release takes time to reach a station. Only S1, and in floods possibly S2, can feel a release
effective within the hour before $\tau^{+}$. The lag windows encode this, so a change in the
announced schedule cannot implausibly affect S4 within the hour.
\end{itemize}

\textbf{(iv) Differences between scheduled and actual releases.}
\begin{itemize}[leftmargin=*,itemsep=1pt]
\item For past steps, the model always sees the \emph{actual} release, plus the release that had
been scheduled for that time as a separate channel ($\mathbf{Q}^{\mathrm{past}}$). It therefore
learns how schedules and actual releases historically differ, and how much to trust them.
\item As soon as a measured release departs from the schedule, the measured value replaces the
schedule for all elapsed steps at the next issuance (at most ten minutes later). The current
deviation, actual minus scheduled, is visible in the input and informs how much the remaining
schedule is trusted.
\item A deviation above an operational threshold triggers an operator alert, a wider interval, and
an optional scenario forecast with the alternative release, which costs only one more forward pass.
\item Training perturbs announced schedules (random offsets and time shifts relative to the actual
releases), so the model does not over-trust announcements.
\end{itemize}

\textbf{Feasibility every ten minutes.} One forward pass handles 7 nodes, 2 gauges, 72 history
steps, about 25 edges with at most 40 admissible lags each, and $d=64$. That is about $10^{7}$
multiply-adds, i.e.\ a few milliseconds on one CPU core. Assembling the inputs takes a handful of
database queries. Even an ensemble of five models plus a few release scenarios stays well under one
second, which leaves almost all of the 10-minute cycle for data collection, quality control and
dissemination. Training is offline (for example monthly, or after a reported channel change, A11)
and takes minutes to hours on a single GPU.
